% Specimen 4: a page loosely modelled on the printed Talmud.
% A short main text sits at the top centre. Two commentaries, in two
% typefaces, wrap around it and widen below it. A band of notes on the
% commentaries runs across the foot.
%
% How it works: each commentary is ONE paragraph with a \parshape. Its
% first n lines are narrow (to clear the main text); the rest are full
% column width. n is computed from the measured height of the main text.
% This composes a single page. It does not break text across pages.
\documentclass[11pt]{article}
\usepackage[paperwidth=7in,paperheight=9in,
            left=0.6in,right=0.6in,top=0.6in,bottom=0.7in]{geometry}
% common.tex -- shared preamble for all layout specimens.
% Each specimen sets its own geometry BEFORE % common.tex -- shared preamble for all layout specimens.
% Each specimen sets its own geometry BEFORE % common.tex -- shared preamble for all layout specimens.
% Each specimen sets its own geometry BEFORE \input{common}.
% Engine: LuaLaTeX (needs fontspec for EB Garamond).

\usepackage{fontspec}
\setmainfont{EB Garamond}[Numbers={OldStyle,Proportional}]
\usepackage{microtype}
\setlength\emergencystretch{1em}
\usepackage{xcolor}
\definecolor{grapeink}{HTML}{5B2A4E}   % one accent color, used sparingly

% Three levels of notes, provided by bigfoot:
%   level 1: arabic numbers   (\footnote)
%   level 2: lowercase letters (\footnoteB)  -- notes on notes
%   level 3: symbols, run-in   (\footnoteC)  -- notes on notes on notes
% bigfoot stacks the series at the page foot in declaration order and
% lets a note of one series contain notes of the next.
\usepackage{bigfoot}
\DeclareNewFootnote{default}
\DeclareNewFootnote{B}[alph]
\DeclareNewFootnote[para]{C}[fnsymbol]
% Letters and symbols restart on every page, so they never run past z or
% past the nine fnsymbol marks.
\usepackage{perpage}
\MakePerPage{footnoteB}
\MakePerPage{footnoteC}

% Semantic note commands used by passage.tex. Each specimen maps them to
% a concrete layout, so the text is written once.
\providecommand\note[1]{\footnote{#1}}
\providecommand\subnote[1]{\footnoteB{#1}}
\providecommand\subsubnote[1]{\footnoteC{#1}}

% A plain, McPhee-like chapter opening: number, title, no ornament.
\newcommand\chapteropening[2]{%
  \clearpage\thispagestyle{empty}%
  \vspace*{0.18\textheight}%
  \begin{center}
    {\large\scshape\color{grapeink} #1}\par\medskip
    {\LARGE #2}\par
  \end{center}
  \vspace{2\baselineskip}\noindent\ignorespaces}

% Specimen label printed in the running foot, so the reader of the
% combined PDF knows which layout a page belongs to.
\newcommand\specimenlabel{}
\usepackage{fancyhdr}
\pagestyle{fancy}\fancyhf{}
\renewcommand\headrulewidth{0pt}
\fancyfoot[C]{\footnotesize\scshape\color{grapeink}\specimenlabel}
\fancypagestyle{empty}{\fancyhf{}\fancyfoot[C]{\footnotesize\scshape\color{grapeink}\specimenlabel}}
.
% Engine: LuaLaTeX (needs fontspec for EB Garamond).

\usepackage{fontspec}
\setmainfont{EB Garamond}[Numbers={OldStyle,Proportional}]
\usepackage{microtype}
\setlength\emergencystretch{1em}
\usepackage{xcolor}
\definecolor{grapeink}{HTML}{5B2A4E}   % one accent color, used sparingly

% Three levels of notes, provided by bigfoot:
%   level 1: arabic numbers   (\footnote)
%   level 2: lowercase letters (\footnoteB)  -- notes on notes
%   level 3: symbols, run-in   (\footnoteC)  -- notes on notes on notes
% bigfoot stacks the series at the page foot in declaration order and
% lets a note of one series contain notes of the next.
\usepackage{bigfoot}
\DeclareNewFootnote{default}
\DeclareNewFootnote{B}[alph]
\DeclareNewFootnote[para]{C}[fnsymbol]
% Letters and symbols restart on every page, so they never run past z or
% past the nine fnsymbol marks.
\usepackage{perpage}
\MakePerPage{footnoteB}
\MakePerPage{footnoteC}

% Semantic note commands used by passage.tex. Each specimen maps them to
% a concrete layout, so the text is written once.
\providecommand\note[1]{\footnote{#1}}
\providecommand\subnote[1]{\footnoteB{#1}}
\providecommand\subsubnote[1]{\footnoteC{#1}}

% A plain, McPhee-like chapter opening: number, title, no ornament.
\newcommand\chapteropening[2]{%
  \clearpage\thispagestyle{empty}%
  \vspace*{0.18\textheight}%
  \begin{center}
    {\large\scshape\color{grapeink} #1}\par\medskip
    {\LARGE #2}\par
  \end{center}
  \vspace{2\baselineskip}\noindent\ignorespaces}

% Specimen label printed in the running foot, so the reader of the
% combined PDF knows which layout a page belongs to.
\newcommand\specimenlabel{}
\usepackage{fancyhdr}
\pagestyle{fancy}\fancyhf{}
\renewcommand\headrulewidth{0pt}
\fancyfoot[C]{\footnotesize\scshape\color{grapeink}\specimenlabel}
\fancypagestyle{empty}{\fancyhf{}\fancyfoot[C]{\footnotesize\scshape\color{grapeink}\specimenlabel}}
.
% Engine: LuaLaTeX (needs fontspec for EB Garamond).

\usepackage{fontspec}
\setmainfont{EB Garamond}[Numbers={OldStyle,Proportional}]
\usepackage{microtype}
\setlength\emergencystretch{1em}
\usepackage{xcolor}
\definecolor{grapeink}{HTML}{5B2A4E}   % one accent color, used sparingly

% Three levels of notes, provided by bigfoot:
%   level 1: arabic numbers   (\footnote)
%   level 2: lowercase letters (\footnoteB)  -- notes on notes
%   level 3: symbols, run-in   (\footnoteC)  -- notes on notes on notes
% bigfoot stacks the series at the page foot in declaration order and
% lets a note of one series contain notes of the next.
\usepackage{bigfoot}
\DeclareNewFootnote{default}
\DeclareNewFootnote{B}[alph]
\DeclareNewFootnote[para]{C}[fnsymbol]
% Letters and symbols restart on every page, so they never run past z or
% past the nine fnsymbol marks.
\usepackage{perpage}
\MakePerPage{footnoteB}
\MakePerPage{footnoteC}

% Semantic note commands used by passage.tex. Each specimen maps them to
% a concrete layout, so the text is written once.
\providecommand\note[1]{\footnote{#1}}
\providecommand\subnote[1]{\footnoteB{#1}}
\providecommand\subsubnote[1]{\footnoteC{#1}}

% A plain, McPhee-like chapter opening: number, title, no ornament.
\newcommand\chapteropening[2]{%
  \clearpage\thispagestyle{empty}%
  \vspace*{0.18\textheight}%
  \begin{center}
    {\large\scshape\color{grapeink} #1}\par\medskip
    {\LARGE #2}\par
  \end{center}
  \vspace{2\baselineskip}\noindent\ignorespaces}

% Specimen label printed in the running foot, so the reader of the
% combined PDF knows which layout a page belongs to.
\newcommand\specimenlabel{}
\usepackage{fancyhdr}
\pagestyle{fancy}\fancyhf{}
\renewcommand\headrulewidth{0pt}
\fancyfoot[C]{\footnotesize\scshape\color{grapeink}\specimenlabel}
\fancypagestyle{empty}{\fancyhf{}\fancyfoot[C]{\footnotesize\scshape\color{grapeink}\specimenlabel}}

\newfontfamily\rightcommentfont{Linux Libertine O}
\renewcommand\specimenlabel{Specimen 4: a Talmud-style page}
\setlength\parindent{0pt}

% ---- page geometry ---------------------------------------------------------
\newlength\mainwidth    \setlength\mainwidth{2.7in}   % central main text
\newlength\colgap       \setlength\colgap{0.2in}       % gap between blocks
\newlength\colwidth     % each commentary column, full width
\setlength\colwidth{\dimexpr(\textwidth-\colgap)/2\relax}
\newlength\narrowwidth  % commentary width beside the main text
\setlength\narrowwidth{\dimexpr(\textwidth-\mainwidth)/2-\colgap\relax}
\newlength\commentskip  \setlength\commentskip{11.6pt}   % commentary leading
\newcommand\commentsize{\fontsize{9.3}{\commentskip}\selectfont}
\newcommand\mainsize{\fontsize{13}{17}\selectfont}
\newcommand\lemma[1]{\textbf{#1}}           % words a comment is keyed to
\newcommand\cmark[1]{\textsuperscript{#1}}  % marks for the notes at the foot

% ---- composition -----------------------------------------------------------
\newsavebox\mainbox
\newcount\narrowlines
% \buildshape{indent}: \shapespec = n narrow lines at that indent, then one
% full-width line (TeX repeats the last line spec for the rest).
\makeatletter
\newcommand\buildshape[1]{% #1 = indent of the narrow lines
  \def\shapespec{}%
  \count@=0
  \loop\ifnum\count@<\narrowlines
    \edef\shapespec{\shapespec #1 \the\narrowwidth}%
    \advance\count@ 1
  \repeat
  \edef\shapespec{\shapespec 0pt \the\colwidth}}
\makeatother
% \talmudpage{main}{left commentary}{right commentary}
\newcommand\talmudpage[3]{%
  \sbox\mainbox{\vtop{\hsize\mainwidth\mainsize
    \emergencystretch=1em #1\par}}%
  % narrow lines = depth of main text below its first baseline, in
  % commentary lines (\numexpr rounds), plus one line for the first
  % baseline and about one line of clearance.
  \narrowlines=\numexpr\dp\mainbox/\commentskip + 2\relax
  \typeout{Talmud page: \the\narrowlines\space narrow commentary lines}
  \noindent\hbox to \textwidth{%
    \rlap{\hskip\dimexpr(\textwidth-\mainwidth)/2\relax\usebox\mainbox}%
    \vtop{\hsize\colwidth\commentsize
      \buildshape{0pt}\parshape\numexpr\narrowlines+1\relax\shapespec
      #2\par}%
    \hfill
    \vtop{\hsize\colwidth\commentsize\rightcommentfont\fontsize{9}{\commentskip}\selectfont
      \buildshape{\dimexpr\colwidth-\narrowwidth\relax}%
      \parshape\numexpr\narrowlines+1\relax\shapespec
      #3\par}}}

\begin{document}
\thispagestyle{empty}
\begin{center}
  {\small\scshape\color{grapeink} Grapes \quad\textperiodcentered\quad
   One \quad\textperiodcentered\quad The Berry}
\end{center}
\vspace{2pt}

\talmudpage
% ---------------------------------------------------------------- main text
{A grape is a berry. Botanists mean something exact by that word: a fleshy
fruit that grows from a single flower with a single ovary, its seeds buried in
the pulp. By this rule the tomato is a berry, and so are the banana and the
eggplant, while the strawberry and the raspberry are not. The grape qualifies
without argument. Nearly all the grapes we eat, dry, and press come from one
species, \emph{Vitis vinifera}, the wine-bearing vine.}
% ------------------------------------------------ left commentary: glosses
{\lemma{A grape.} The single berry, not the bunch. The English word comes
from French \emph{grappe}, the bunch, so that the name of the cluster became
the name of the berry. \lemma{Is a berry.} In the botanist's sense, not the
grocer's. The grocer calls a berry anything small, soft, and sweet that can
be eaten without peeling. \lemma{Something exact.} The definition here is
the textbook one. Botanists argue about the hard cases. \lemma{A single
flower with a single ovary.} This separates a true berry from an aggregate
fruit such as the raspberry, which grows from one flower with many ovaries,
and from a multiple fruit such as the mulberry or the pineapple, which grows
from many flowers fused together. \lemma{Seeds buried in the pulp.} Compare
the stone fruits, the cherry and the plum, whose single seed is locked in a
hard pit. \lemma{The tomato.} The Supreme Court of the United States once
ruled that the tomato is a vegetable, in the case of Nix v.\ Hedden, decided
in 1893.\cmark{*} The court did not dispute the botany. It held that in
ordinary speech, and therefore in the law, tomatoes are vegetables.
\lemma{The banana.} The seeds of wild bananas are large and hard. The small
dark specks in a cultivated banana are the remains of seeds that no longer
develop.\cmark{\dag} \lemma{The eggplant.} Like the tomato, a member of the
nightshade family. \lemma{The strawberry.} Its flesh grows from the
receptacle, the part of the flower that holds the ovaries, and the specks on
its skin are the true fruits. \lemma{Qualifies without argument.} Which is
more than can be said of most things in this book. \lemma{Vitis vinifera.}
\emph{Vitis} is the Latin for vine, and \emph{vinifera} means wine-bearing.
The rest of the genus, several dozen species, grows wild across the temperate
Northern Hemisphere. Many of them are American, and this fact will matter a
great deal later.}
% --------------------------------------- right commentary: on this page
{\lemma{On this page.} The page imitates, loosely, the printed Talmud, where a
short passage of main text sits in the centre and the commentaries wrap
around it. The main text is set first. Each commentary is then set as one
long paragraph whose first lines are narrowed to leave room for it; below the
main text the commentaries widen to fill the half page. \lemma{Two
commentaries.} Here the left column glosses single words, and this column
makes larger remarks. A book about grapes might give the left column to
botany and the right to history, or the left to fact and the right to
digression. \lemma{Keyed by lemma.} The commentaries use no note numbers.
Each comment opens with the words it comments on, in bold, and the reader
finds the place by matching words. The main text stays clean. The reader pays
for this with a little searching. \lemma{Two typefaces.} The left column is
set in EB Garamond and this one in Linux Libertine, so the eye can tell them
apart without labels. Printed Talmuds do something similar: the main text and
the commentaries are set in different Hebrew scripts.\cmark{\ddag}
\lemma{The band at the foot.} It holds a third level: notes on the
commentaries, marked with symbols. \lemma{What this costs.} Each page must be
composed as a unit. The main text for a page has to be chosen by hand, and the
commentaries must be long enough to fill the page but not longer. For a whole
book this would need either a great deal of hand work or a small program that
measures the text and cuts it into pages.}

\vfill
\noindent\rule{\textwidth}{0.3pt}\par\vspace{3pt}
{\fontsize{8}{9.6}\selectfont
\cmark{*}\,The tariff at issue, from 1883, taxed imported vegetables but not
fruit. \quad
\cmark{\dag}\,Most bananas sold in Europe and North America belong to one
group of cultivars, the Cavendish. \quad
\cmark{\ddag}\,The commentaries are traditionally set in what is called Rashi
script, a semi-cursive Hebrew typeface. The name is a later convention; Rashi
himself did not write in it.\par}
\end{document}
